\section{Бинарные операции и их свойства}
\label{sec:BinaryOperations}
\tikzsetfigurename{LinearAlgebra_BinaryOp_}

Введём понятие \emph{бинарной операции} и рассмотрим некоторые её свойства, такие как \emph{коммутативность} и \emph{ассоциативность}.

\begin{definition}[Функция]
    \emph{Функцией} будем называть бинарное отношение на двух множествах $S$ и $T$, такое, что каждому элементу из $S$ сопоставляется ровно один элемент из $T$.
    Запись $f: S \to T$ как раз и обозначает, что функция $f$ сопоставляет элементы из $S$ элементам из $T$.
\end{definition}

\begin{definition}[Домен функции]
    Для функции $f: S \to T$, множество $S$ называется \emph{областью определения функции} или \emph{доменом функции}.
\end{definition}

\begin{definition}[Кодомен функции]
    Для функции $f: S \to T$, множество $T$ называется \emph{областью значений функции} или \emph{кодоменом функции}.
\end{definition}

\begin{definition}[Инъекция]
    Функция $f: S \to T$ называется \emph{инъективной} или \emph{инъекцией}, если она отображает разные элементы из $S$ в разные элементы из $T$.
    Иными словами, для любых $s_1, s_2 \in S$ из $s_1 \neq s_2$ следует $f(s_1) \neq f(s_2)$.
\end{definition}

\begin{definition}[Сюръекция]
    Функция $f: S \to T$ называется \emph{сюръективной} или \emph{сюръекцией}, если каждый элемент из $T$ является образом хотя бы одного элемента из $S$.
    Иными словами, для любого $t \in T$ существует $s \in S$ такой, что $f(s) = t$.
\end{definition}

\begin{definition}[Биекция]
    Функция $f: S \to T$ называется \emph{биективной} или \emph{биекцией}, если она является одновременно инъекцией и сюръекцией.
    Иными словами, биекция~--- это взаимно-однозначное соответствие между множествами $S$ и $T$.
\end{definition}

\begin{definition}[Двухместная функция]
    Функцию, принимающую два аргумента, $f: R \times S \to T$ будем называть \emph{двухместной} или \emph{функцией арности два}.
    Для записи таких функций будем использовать типичную нотацию: $t = f(r, s)$.
\end{definition}

\begin{definition}[Бинарная операция]
    \emph{Бинарная операция}~--- это двухместная функция, от которой дополнительно требуется, чтобы оба аргумента и результат лежали в одном и том же множестве: $f: S \times S \to S$.
    В таком случае говорят, что бинарная операция определена на некотором множестве $S$. Для обозначения произвольной бинарной операции будем использовать символ $\circ$ и пользоваться инфиксной нотацией для записи: $s_3 = s_1 \circ s_2$.
\end{definition}

\begin{definition}[Внешняя бинарная операция]
    \emph{Внешняя бинарная операция}~--- это бинарная операция, у которой аргументы лежат в разных множествах, при этом результат~--- в одном из этих множеств.
    Иными словами $\circ: R \times S \to S$, где $R$ может не совпадать с $S$~--- внешняя бинарная операция.
\end{definition}

Следует помнить, что как функции, так и бинарные операции, могут быть частично определёнными (частичные функции, частичные бинарные операции).
Типичным примером частично определённой бинарной операции является деление на целых числах: она не определена, если второй аргумент равен нулю.

Бинарные операции могут обладать некоторыми дополнительными свойствами, такими как \emph{коммутативность} или \emph{ассоциативность}, позволяющими преобразовывать выражения, составленные с использованием этих операций.

\begin{definition}[Коммутативная операция]
    Бинарная операция $\circ : S \times S \to S$ называется \emph{коммутативной}, если для любых $s_1, s_2 \in S$ верно, что $s_1 \circ s_2 = s_2 \circ s_1$.
\end{definition}

\begin{example}
    Рассмотрим несколько примеров коммутативных и некоммутативных операций.
    \begin{itemize}
        \item Операция сложения на целых числах является коммутативной: известный ещё со школы перестановочный закон сложения.
        \item Операция умножения на целых числах является коммутативной: известный ещё со школы перестановочный закон умножения.
        \item Операция конкатенации на строках\sidenote{Хотя в языках программирования традиционно строки принято <<складывать>>, использование $\cdot$ позволяет нам естественным образом пользоваться математической традицией опускать знак умножения в символьных записях.} $\cdot$ не является коммутативной:
              \["ab" \cdot "c"  = "abc" \neq "cab" = "c" \cdot "ab".\]
        \item Операция умножения матриц (над целыми числами) $\cdot$ не является коммутативной:
              \[\begin{pmatrix}
                      1 & 1 \\
                      0 & 0
                  \end{pmatrix}
                  \cdot
                  \begin{pmatrix}
                      0 & 0 \\
                      1 & 1
                  \end{pmatrix}
                  =
                  \begin{pmatrix}
                      1 & 1 \\
                      0 & 0
                  \end{pmatrix}
                  \neq
                  \begin{pmatrix}
                      0 & 0 \\
                      1 & 1
                  \end{pmatrix}
                  =
                  \begin{pmatrix}
                      0 & 0 \\
                      1 & 1
                  \end{pmatrix}
                  \cdot
                  \begin{pmatrix}
                      1 & 1 \\
                      0 & 0
                  \end{pmatrix}
                  .\]
    \end{itemize}
\end{example}

\begin{definition}[Ассоциативная бинарная операция]
    Бинарная операция $\circ: S \times S \to S$ называется \emph{ассоциативной}, если для любых $s_1, s_2, s_3 \in S$ верно, что $(s_1 \circ s_2) \circ s_3 = s_1 \circ (s_2 \circ s_3)$.
    Иными словами, для ассоциативной операции результат вычислений не зависит от порядка применения операций.
\end{definition}

\begin{example} Рассмотрим несколько примеров ассоциативных и неассоциативных операций.
    \begin{itemize}
        \item Операция сложения на целых числах является ассоциативной.
        \item Операция умножения на целых числах является ассоциативной.
        \item Операция конкатенации на строках $\cdot$ является ассоциативной:
              \[("a" \cdot "b") \cdot "c"  = "a" \cdot ("b" \cdot "c") = "abc".\]
        \item Операция возведения в степень (над целыми числами) $\hat{\mkern6mu}$ не является ассоциативной:
              \[(2\hat{\mkern6mu}2)\hat{\mkern6mu}3 = 4 \hat{\mkern6mu} 3 = 64 \neq 256 = 2 \hat{\mkern6mu} 8 = 2\hat{\mkern6mu}(2\hat{\mkern6mu}3).\]
    \end{itemize}
\end{example}

\begin{definition}[Дистрибутивная бинарная операция]
    Говорят, что бинарная операция $\otimes: S \times S \to S$ является \emph{дистрибутивной} относительно бинарной операции $\oplus: S \times S \to S$, если
    \begin{enumerate}
        \item Для любых $s_1, s_2, s_3 \in S$, $s_1 \otimes (s_2 \oplus s_3) = (s_1 \otimes s_2) \oplus (s_1 \otimes s_3)$ (дистрибутивность слева).
        \item Для любых $s_1, s_2, s_3 \in S$, $(s_2 \oplus s_3) \otimes s_1 = (s_2 \otimes s_1) \oplus (s_3 \otimes s_1)$ (дистрибутивность справа).
    \end{enumerate}
    Если операция $\otimes$ является коммутативной, то дистрибутивность слева и справа равносильны.
\end{definition}

\begin{example}
    Рассмотрим несколько примеров дистрибутивных операций.
    \begin{itemize}
        \item Умножение целых чисел дистрибутивно относительно сложения и вычитания: классический \emph{распределительный закон}, знакомый всем со школы.
        \item Операция деления (допустим, на действительных числах) не коммутативна.
              При этом, она дистрибутивна справа относительно сложения и вычитания, но не дистрибутивна слева%
              \sidenote{
                  Здесь может быть уместно вспомнить правила сложения дробей.
                  Дроби с общим знаменателем складывать проще как раз из-за дистрибутивности справа.}.
              Так, $(a + b) / c = (a / c) + (b / c)$, но $c / (a + b) \neq (c / a) + (c / b)$.
    \end{itemize}
\end{example}

\begin{definition}[Идемпотентная бинарная операция]
    Бинарная операция $\circ: S \times S \to S$ называется \emph{идемпотентной}, если для любого $s \in S$ верно, что $s \circ s = s$.
\end{definition}

\begin{example}
    Рассмотрим несколько примеров идемпотентных операций.
    \begin{itemize}
        \item Операция объединения множеств $\cup$ является идемпотентной: для любого множества $S$ верно, что $S \cup S = S$.
        \item Операция сложения на целых числах не является идемпотентной.
        \item Операции \enquote{логическое и} $\land$ и \enquote{логическое или} $\lor$ являются идемпотентными.
        \item Операция \enquote{исключающее или} (\textsf{XOR}) не является идемпотентной.
    \end{itemize}
\end{example}

\begin{definition}[Нейтральный элемент]
    Пусть есть коммутативная бинарная операция $\circ$ на множестве $S$.
    Говорят, что $e \in S$ является \emph{нейтральным элементом} по операции $\circ$, если для любого $s \in S$ верно, что $e \circ s = s \circ e = s$.
    Если бинарная операция не является коммутативной, то можно определить \emph{нейтральный слева} и \emph{нейтральный справа} элементы по аналогии.
\end{definition}

